% Lecture 11 - Page 2 (13-Dec-1994).
% Fragment for % Lecture 11 - Page 2 (13-Dec-1994).
% Fragment for % Lecture 11 - Page 2 (13-Dec-1994).
% Fragment for % Lecture 11 - Page 2 (13-Dec-1994).
% Fragment for \input{courses/pages/curs11-pagina2-en}.
% No preamble or forced page break.
% The manuscript's sequence is preserved; uncertain readings are noted.

\subsection*{Lecture 11 - Page 2}

\noindent I claim that $|A_0|=1 \Longleftrightarrow
(\forall)\,Q\in A_0$ we have $Q=P$.

\medskip
\[
P\subseteq N_G(Q)=\{a\in G\mid aQa^{-1}=Q\}\supseteq Q.
\]
\[
N_G(Q)\leq G \quad\Longrightarrow\quad |N_G(Q)|\mid p^t m.
\]
\[
|P|=p^t\mid |N_G(Q)|.
\]
\[
\Longrightarrow\quad |N_G(Q)|=p^t m_1,
\qquad p\nmid m_1,\quad m_1\mid m.
\]

\noindent $\Longrightarrow$ $P,Q$ are Sylow $p$-subgroups of $N_G(Q)$.

\medskip
\noindent By Sylow II $\Longrightarrow$ there exists $a\in N_G(Q)$ such that
\[
aQa^{-1}=P.
\]
\noindent $a\in N_G(Q) \Longrightarrow aQa^{-1}=Q
\Longrightarrow P=Q \Longrightarrow |A_0|=1$.

\medskip
\noindent\underline{We deduce:}
\[
|\operatorname{Syl}_p(G)|\equiv1\pmod p.
\]
% In the handwritten concluding line the cardinality bars are absent
% or indistinct: it appears as "Syl_p(G) = 1 (mod p)".

\bigskip
\noindent\underline{Application:}
Let $G$ be a finite group, $|G|=pq$, $p>q$, $q\nmid p-1$,
where $p,q$ are primes. \textit{(35)}
\[
\text{Then}\qquad G\cong (\mathbb{Z}_{pq},+).
\]
% The handwritten parenthetical "(35)" is not unambiguously identifiable;
% it is retained as an annotation, without assigning a meaning to it.

\noindent $A=\operatorname{Syl}_p(G),\qquad B=\operatorname{Syl}_q(G)$.

\medskip
\noindent By Cauchy's theorem $\Longrightarrow$ there exist $a,b\in G$
with $\operatorname{ord}(a)=p$, $\operatorname{ord}(b)=q$.

\[
\begin{aligned}
P&=\langle a\rangle\leq G, & P&\in A,\\
Q&=\langle b\rangle\leq G, & Q&\in B.
\end{aligned}
\]

\noindent By Sylow III:
$|A|\mid pq$ and $|B|\mid pq$, hence
\[
\left\{
\begin{aligned}
|A|&\in\{1,p,q,pq\},\\
|B|&\in\{1,p,q,pq\}.
\end{aligned}
\right.
\]

\noindent I shall show that $|A|=|B|=1$.

\medskip
\noindent Consider $[a,b]=aba^{-1}b^{-1}$.

\medskip
\noindent $P\trianglelefteq G$,\quad
$xPx^{-1}\in\operatorname{Syl}_p(G)$, $(\forall)\,x\in G$.
% The manuscript page ends here. The argument showing |A|=|B|=1
% and the conclusion of the application are not present on this page.
.
% No preamble or forced page break.
% The manuscript's sequence is preserved; uncertain readings are noted.

\subsection*{Lecture 11 - Page 2}

\noindent I claim that $|A_0|=1 \Longleftrightarrow
(\forall)\,Q\in A_0$ we have $Q=P$.

\medskip
\[
P\subseteq N_G(Q)=\{a\in G\mid aQa^{-1}=Q\}\supseteq Q.
\]
\[
N_G(Q)\leq G \quad\Longrightarrow\quad |N_G(Q)|\mid p^t m.
\]
\[
|P|=p^t\mid |N_G(Q)|.
\]
\[
\Longrightarrow\quad |N_G(Q)|=p^t m_1,
\qquad p\nmid m_1,\quad m_1\mid m.
\]

\noindent $\Longrightarrow$ $P,Q$ are Sylow $p$-subgroups of $N_G(Q)$.

\medskip
\noindent By Sylow II $\Longrightarrow$ there exists $a\in N_G(Q)$ such that
\[
aQa^{-1}=P.
\]
\noindent $a\in N_G(Q) \Longrightarrow aQa^{-1}=Q
\Longrightarrow P=Q \Longrightarrow |A_0|=1$.

\medskip
\noindent\underline{We deduce:}
\[
|\operatorname{Syl}_p(G)|\equiv1\pmod p.
\]
% In the handwritten concluding line the cardinality bars are absent
% or indistinct: it appears as "Syl_p(G) = 1 (mod p)".

\bigskip
\noindent\underline{Application:}
Let $G$ be a finite group, $|G|=pq$, $p>q$, $q\nmid p-1$,
where $p,q$ are primes. \textit{(35)}
\[
\text{Then}\qquad G\cong (\mathbb{Z}_{pq},+).
\]
% The handwritten parenthetical "(35)" is not unambiguously identifiable;
% it is retained as an annotation, without assigning a meaning to it.

\noindent $A=\operatorname{Syl}_p(G),\qquad B=\operatorname{Syl}_q(G)$.

\medskip
\noindent By Cauchy's theorem $\Longrightarrow$ there exist $a,b\in G$
with $\operatorname{ord}(a)=p$, $\operatorname{ord}(b)=q$.

\[
\begin{aligned}
P&=\langle a\rangle\leq G, & P&\in A,\\
Q&=\langle b\rangle\leq G, & Q&\in B.
\end{aligned}
\]

\noindent By Sylow III:
$|A|\mid pq$ and $|B|\mid pq$, hence
\[
\left\{
\begin{aligned}
|A|&\in\{1,p,q,pq\},\\
|B|&\in\{1,p,q,pq\}.
\end{aligned}
\right.
\]

\noindent I shall show that $|A|=|B|=1$.

\medskip
\noindent Consider $[a,b]=aba^{-1}b^{-1}$.

\medskip
\noindent $P\trianglelefteq G$,\quad
$xPx^{-1}\in\operatorname{Syl}_p(G)$, $(\forall)\,x\in G$.
% The manuscript page ends here. The argument showing |A|=|B|=1
% and the conclusion of the application are not present on this page.
.
% No preamble or forced page break.
% The manuscript's sequence is preserved; uncertain readings are noted.

\subsection*{Lecture 11 - Page 2}

\noindent I claim that $|A_0|=1 \Longleftrightarrow
(\forall)\,Q\in A_0$ we have $Q=P$.

\medskip
\[
P\subseteq N_G(Q)=\{a\in G\mid aQa^{-1}=Q\}\supseteq Q.
\]
\[
N_G(Q)\leq G \quad\Longrightarrow\quad |N_G(Q)|\mid p^t m.
\]
\[
|P|=p^t\mid |N_G(Q)|.
\]
\[
\Longrightarrow\quad |N_G(Q)|=p^t m_1,
\qquad p\nmid m_1,\quad m_1\mid m.
\]

\noindent $\Longrightarrow$ $P,Q$ are Sylow $p$-subgroups of $N_G(Q)$.

\medskip
\noindent By Sylow II $\Longrightarrow$ there exists $a\in N_G(Q)$ such that
\[
aQa^{-1}=P.
\]
\noindent $a\in N_G(Q) \Longrightarrow aQa^{-1}=Q
\Longrightarrow P=Q \Longrightarrow |A_0|=1$.

\medskip
\noindent\underline{We deduce:}
\[
|\operatorname{Syl}_p(G)|\equiv1\pmod p.
\]
% In the handwritten concluding line the cardinality bars are absent
% or indistinct: it appears as "Syl_p(G) = 1 (mod p)".

\bigskip
\noindent\underline{Application:}
Let $G$ be a finite group, $|G|=pq$, $p>q$, $q\nmid p-1$,
where $p,q$ are primes. \textit{(35)}
\[
\text{Then}\qquad G\cong (\mathbb{Z}_{pq},+).
\]
% The handwritten parenthetical "(35)" is not unambiguously identifiable;
% it is retained as an annotation, without assigning a meaning to it.

\noindent $A=\operatorname{Syl}_p(G),\qquad B=\operatorname{Syl}_q(G)$.

\medskip
\noindent By Cauchy's theorem $\Longrightarrow$ there exist $a,b\in G$
with $\operatorname{ord}(a)=p$, $\operatorname{ord}(b)=q$.

\[
\begin{aligned}
P&=\langle a\rangle\leq G, & P&\in A,\\
Q&=\langle b\rangle\leq G, & Q&\in B.
\end{aligned}
\]

\noindent By Sylow III:
$|A|\mid pq$ and $|B|\mid pq$, hence
\[
\left\{
\begin{aligned}
|A|&\in\{1,p,q,pq\},\\
|B|&\in\{1,p,q,pq\}.
\end{aligned}
\right.
\]

\noindent I shall show that $|A|=|B|=1$.

\medskip
\noindent Consider $[a,b]=aba^{-1}b^{-1}$.

\medskip
\noindent $P\trianglelefteq G$,\quad
$xPx^{-1}\in\operatorname{Syl}_p(G)$, $(\forall)\,x\in G$.
% The manuscript page ends here. The argument showing |A|=|B|=1
% and the conclusion of the application are not present on this page.
.
% No preamble or forced page break.
% The manuscript's sequence is preserved; uncertain readings are noted.

\subsection*{Lecture 11 - Page 2}

\noindent I claim that $|A_0|=1 \Longleftrightarrow
(\forall)\,Q\in A_0$ we have $Q=P$.

\medskip
\[
P\subseteq N_G(Q)=\{a\in G\mid aQa^{-1}=Q\}\supseteq Q.
\]
\[
N_G(Q)\leq G \quad\Longrightarrow\quad |N_G(Q)|\mid p^t m.
\]
\[
|P|=p^t\mid |N_G(Q)|.
\]
\[
\Longrightarrow\quad |N_G(Q)|=p^t m_1,
\qquad p\nmid m_1,\quad m_1\mid m.
\]

\noindent $\Longrightarrow$ $P,Q$ are Sylow $p$-subgroups of $N_G(Q)$.

\medskip
\noindent By Sylow II $\Longrightarrow$ there exists $a\in N_G(Q)$ such that
\[
aQa^{-1}=P.
\]
\noindent $a\in N_G(Q) \Longrightarrow aQa^{-1}=Q
\Longrightarrow P=Q \Longrightarrow |A_0|=1$.

\medskip
\noindent\underline{We deduce:}
\[
|\operatorname{Syl}_p(G)|\equiv1\pmod p.
\]
% In the handwritten concluding line the cardinality bars are absent
% or indistinct: it appears as "Syl_p(G) = 1 (mod p)".

\bigskip
\noindent\underline{Application:}
Let $G$ be a finite group, $|G|=pq$, $p>q$, $q\nmid p-1$,
where $p,q$ are primes. \textit{(35)}
\[
\text{Then}\qquad G\cong (\mathbb{Z}_{pq},+).
\]
% The handwritten parenthetical "(35)" is not unambiguously identifiable;
% it is retained as an annotation, without assigning a meaning to it.

\noindent $A=\operatorname{Syl}_p(G),\qquad B=\operatorname{Syl}_q(G)$.

\medskip
\noindent By Cauchy's theorem $\Longrightarrow$ there exist $a,b\in G$
with $\operatorname{ord}(a)=p$, $\operatorname{ord}(b)=q$.

\[
\begin{aligned}
P&=\langle a\rangle\leq G, & P&\in A,\\
Q&=\langle b\rangle\leq G, & Q&\in B.
\end{aligned}
\]

\noindent By Sylow III:
$|A|\mid pq$ and $|B|\mid pq$, hence
\[
\left\{
\begin{aligned}
|A|&\in\{1,p,q,pq\},\\
|B|&\in\{1,p,q,pq\}.
\end{aligned}
\right.
\]

\noindent I shall show that $|A|=|B|=1$.

\medskip
\noindent Consider $[a,b]=aba^{-1}b^{-1}$.

\medskip
\noindent $P\trianglelefteq G$,\quad
$xPx^{-1}\in\operatorname{Syl}_p(G)$, $(\forall)\,x\in G$.
% The manuscript page ends here. The argument showing |A|=|B|=1
% and the conclusion of the application are not present on this page.
